% Chapter 2 of Optimizers in the Wild: sign descent and the l-infinity geometry of Adam, signSGD and Lion.

\casechapter{Sign Descent}{The $\ell_\infty$ geometry of Adam, signSGD and Lion}{signSGD, Lion, Adam without memory}{ch:sign}
\begin{refsection}[chapters/02-sign.bib]

\begin{chaptersummary}
Adam without its moving averages, signSGD and Lion move every coordinate by the same amount: they take the step $-\eta\,\sign(\cdot)$, a step of fixed length in the max norm. We prove that this step gains $\eta\norm{\nabla\Loss}_1$ against a curvature cost $\frac12\eta^2\norm{\vec L}_1$, that the trade is exact on separable quadratics, and that a rotation of the coordinates can raise the cost by the factor $d\lmax/\tr H$, which gradient descent does not notice. With minibatches the sign follows the median of the gradient rather than its mean; the resulting bias $\norm{\vec\sigma}_1/\sqrt B$ is sharp, a rare class realizes it, and symmetric noise avoids it. Adam's step is at most $C(\beta_1,\beta_2)$ times the learning rate in each coordinate, $7.27$ at the defaults; the idealization drops the memory that makes Adam work, and it cannot say why the coordinates of real networks are good ones.
\end{chaptersummary}

\section{In the wild: one step length for every coordinate}\label{sec:sign:wild}

Adam divides a moving average of gradients by the square root of a moving average of squared gradients \parencite{KingmaBa2015}. The division removes the scale of each coordinate: multiply the gradient of one parameter by a thousand, and Adam's step in that parameter does not change. Its learning rate is therefore not a multiplier of the gradient but a length, roughly the distance each parameter moves in one step. Two other optimizers make the same choice exactly. \term{SignSGD}\index{signSGD} updates with the sign of a stochastic gradient; in data-parallel training each worker then sends one bit per coordinate, and a server that takes the majority vote of the signs sends one bit back \parencite{BernsteinEtAl2018,BernsteinEtAl2019}. The idea goes back to Rprop, which used only the signs of a full-batch gradient,\sidenote{Rprop keeps a step size for each weight, grows it while the weight's partial derivative keeps its sign and shrinks it when the sign flips; the size of the gradient never enters the update \parencite{RiedmillerBraun1993}. RMSProp was introduced in a lecture course as a minibatch version of Rprop \parencite{TielemanHinton2012}.} and to 1-bit SGD, which quantized each gradient coordinate to one bit for distributed training and carried the quantization error over to the next step \parencite{SeideEtAl2014}. Lion, found by a symbolic search over update rules, updates with the sign of an interpolation between the gradient and a momentum, and keeps one number of state per parameter where Adam keeps two \parencite{ChenEtAl2023}.

\begin{practice}[Adam, \textcite{KingmaBa2015}; GPT-3, \textcite{BrownEtAl2020}; Lion, \textcite{ChenEtAl2023}; signSGD, \textcite{BernsteinEtAl2019}]
Adam with $\beta_1=0.9$, $\beta_2=0.999$, $\epsilon=10^{-8}$, the published defaults, and $\beta_2=0.95$ in large language models since GPT-3. Lion with $\beta_1=0.9$, $\beta_2=0.99$, a learning rate $3$ to $10$ times smaller than AdamW's and a weight decay $3$ to $10$ times larger, so that their product, the strength of the decay, stays about the same. In distributed training, send $\sign(g)$, one bit per coordinate, and aggregate by majority vote.
\end{practice}

Three observations tie these optimizers together.

\begin{observation}[Sign descent with momentum behaves like Adam; \textcite{KunstnerEtAl2023}, \textcite{ZhaoEtAl2025}]\label[observation]{obs:sign:signum}
On transformer language models the gap between Adam and stochastic gradient descent does not shrink as the batch grows to the full training set: Adam improves with the batch size and gradient descent much less, and with large batches Adam behaves like sign descent with momentum. In a comparison of optimizers for autoregressive language models across sizes and hyperparameters, signed momentum (Signum) recovered both the performance of Adam and its stability across hyperparameters.
\end{observation}

\Textcite{KunstnerEtAl2023} concluded that heavy-tailed noise, proposed earlier as the reason for the gap, is not its main cause, and \textcite{ZhaoEtAl2025} found Adam, Adafactor, Lion and Sophia comparable, with only stochastic gradient descent behind. Both studies are empirical and at moderate scale, and neither says why the sign helps.

\begin{observation}[Adam is a sign with a variance-dependent length; \textcite{BallesHennig2018}]\label[observation]{obs:sign:dissect}
Adam's update of each coordinate can be read as the sign of the momentum times a factor that decreases with the relative variance of that coordinate's gradient. Studied separately, the sign is the aspect responsible for adverse effects on generalization that had been reported for adaptive methods.
\end{observation}

\begin{observation}[Rotating the parameters slows Adam; \textcite{XieMohamadiLi2025}]\label[observation]{obs:sign:rotation}
Training GPT-2 with Adam after a rotation of the parameter space is markedly slower than in the original coordinates, while stochastic gradient descent, whose trajectory rotates with the coordinates, is unaffected.
\end{observation}

The three point at one object, the step $-\eta\,\sign(g)$, and one question: what does it buy, compared with $-\eta g$, and what does it cost? The answer depends on the norm in which steps are measured, on the coordinates, on the shape of the noise, and on memory.

\section{The idealization: steepest descent in the max norm}\label{sec:sign:ideal}

\Cref{ch:steepest} turned a gradient into a step with a norm: the normalized steepest-descent step is $\eta\lmo(g)$ with $\lmo(g)\in\argmin_{\norm u\le1}\ip{g,u}$. For the max norm the dual norm is $\norm{g}_1$ (\cref{lem:toolkit:dualnorm}), and $\ip{g,u}$ is minimized over the cube $\{\norm{u}_\infty\le1\}$ by $u_i=-\sign(g_i)$ where $g_i\ne0$ and by every $u_i\in[-1,1]$ where $g_i=0$; the convention $\sign0=0$ picks the shortest minimizer. Normalized steepest descent in $\ell_\infty$ is sign descent.

\begin{definition}[Sign descent and signSGD]\label[definition]{def:sign:signdescent}
Given step sizes $\eta_t>0$, \term{sign descent}\index{sign descent} is $\theta_{t+1}=\theta_t-\eta_t\sign(\nabla\Loss(\theta_t))$, and signSGD is $\theta_{t+1}=\theta_t-\eta_t\sign(g_t)$ for a stochastic gradient $g_t$, with the sign taken coordinatewise. When no coordinate of the gradient vanishes, the step has length $\eta_t$ in $\ell_\infty$ and in the RMS norm, and $\eta_t\sqrt d$ in $\ell_2$.
\end{definition}

The learning rate of the practice sections and the step size of the theory are the same number $\eta$; for a sign method it is the distance that every coordinate moves. The curvature bound that suits the max norm is the one with which \textcite{BernsteinEtAl2018} posed the problem.

\begin{assumption}[Coordinatewise smoothness]\label[assumption]{ass:sign:smooth}
The loss $\Loss\colon\R^d\to\R$ is differentiable and bounded below by $\Loss^\star$, and there are $L_1,\dots,L_d\ge0$ such that for all $\theta,u\in\R^d$
\[
  \Loss(\theta+u)\le\Loss(\theta)+\ip{\nabla\Loss(\theta),u}+\frac12\sum_{i=1}^dL_iu_i^2 .
\]
We write $\vec L=(L_1,\dots,L_d)$, $\norm{\vec L}_1=\sum_iL_i$ and $L_{\max}=\max_iL_i$.
\end{assumption}

For a twice-differentiable loss, \term{coordinatewise smoothness}\index{coordinatewise smoothness} says that $\nabla^2\Loss(\theta)\preceq\diag(\vec L)$ at every $\theta$. Since $\frac12\sum_iL_iu_i^2$ is at most $\frac12\norm{\vec L}_1\norm{u}_\infty^2$ and at most $\frac12L_{\max}\norm{u}_2^2$, it gives the quadratic upper bounds of \cref{ch:steepest} (\cref{lem:toolkit:descent}) in two geometries at once: the constant $\norm{\vec L}_1$ in $\ell_\infty$ and $L_{\max}$ in $\ell_2$.

\begin{remark}[Coordinatewise and $\ell_\infty$ smoothness]\label[remark]{rem:sign:sdp}
For a quadratic with Hessian $H\succeq0$ the admissible $\vec L$ are those with $\diag(\vec L)\succeq H$, and each has $\norm{\vec L}_1=s^\top\diag(\vec L)s\ge s^\top Hs$ for every sign vector $s$. Hence $\norm{\vec L}_1\ge\norm{H}_{\infty\to1}=\max_{s\in\{\pm1\}^d}s^\top Hs$, the exact $\ell_\infty$ constant of \cref{ch:steepest}, with equality for diagonal $H$ and $\vec L=\diag(H)$; conversely the best $\norm{\vec L}_1$ is at most $\frac\pi2\norm{H}_{\infty\to1}$ (\cref{ex:sign:sdp}, after \textcite{Nesterov1998}).
\end{remark}

\begin{assumption}[Minibatch noise]\label[assumption]{ass:sign:noise}
The stochastic gradient is $g_t=\nabla\Loss(\theta_t)+\zeta_t$ where, conditionally on $\theta_0,g_0,\dots,g_{t-1}$, each coordinate $\zeta_{t,i}$ has mean zero and variance at most $\sigma_i^2/B$. We write $\vec\sigma=(\sigma_1,\dots,\sigma_d)$.
\end{assumption}

A batch of $B$ independent examples whose per-example gradients have coordinate variances at most $\sigma_i^2$ satisfies the assumption. Adam enters with the conventions of the book: $m_{-1}=v_{-1}=0$, $m_t=\beta_1m_{t-1}+(1-\beta_1)g_t$, $v_t=\beta_2v_{t-1}+(1-\beta_2)g_t^2$ with squares taken coordinatewise, $\hat m_t=m_t/(1-\beta_1^{t+1})$, $\hat v_t=v_t/(1-\beta_2^{t+1})$, and $\theta_{t+1}=\theta_t-\eta\,\hat m_t/(\sqrt{\hat v_t}+\epsilon)$.

\begin{dictionary}
the step $-\eta\sign(g)$, of length $\eta$ in $\ell_\infty$ & an update that moves every parameter by the learning rate\\
the dual norm $\norm{\nabla\Loss}_1$ & the progress per unit of learning rate, summed over parameters\\
$\vec L$ and $\norm{\vec L}_1$ & the curvature that each parameter sees along its own axis, and its sum\\
$\norm{g}_1^2/\norm{g}_2^2$, between $1$ and $d$ & the number of parameters that receive a gradient of comparable size\\
$\tr H/\lmax$ & the effective number of stiff directions\\
an orthogonal matrix $Q$ & a change of basis of parameter space, such as mixing the weights of a layer\\
$\sigma_i/\sqrt B$; the median of $g_{t,i}$ & the minibatch noise of a parameter; what a sign update follows\\
$C(\beta_1,\beta_2)$ & Adam's largest step of a parameter, in units of the learning rate
\end{dictionary}

\begin{fieldguide}{signSGD and Signum}{\textcite{BernsteinEtAl2018}}
\fgrow{update}{$\theta_{t+1}=\theta_t-\eta\,\sign(g_t)$; Signum: $\theta_{t+1}=\theta_t-\eta\,\sign(m_t)$, $m_t=\beta m_{t-1}+(1-\beta)g_t$}
\fgrow{norm}{$\ell_\infty$; $\lmo(g)=-\sign(g)$}
\fgrow{preconditioner}{none (equivalently $\diag\abs{g_t}^{-1}$, rebuilt at every step)}
\fgrow{averages}{none (signSGD); EMA of $g$ (Signum)}
\fgrow{state}{none (signSGD); one number per parameter (Signum)}
\fgrow{variants}{majority vote over $M$ workers, $\theta_{t+1}=\theta_t-\eta\,\sign\bigl(\sum_m\sign(g^{(m)}_t)\bigr)$ \parencite{BernsteinEtAl2019}; error feedback \parencite{SeideEtAl2014,KarimireddyEtAl2019}}
\fgrow{caution}{without momentum the step follows the median of the minibatch gradient, not its mean (\cref{prop:sign:rareclass})}
\end{fieldguide}

\begin{fieldguide}{Lion}{\textcite{ChenEtAl2023}}
\fgrow{update}{$c_t=\beta_1m_{t-1}+(1-\beta_1)g_t$, \ $m_t=\beta_2m_{t-1}+(1-\beta_2)g_t$,\\ $\theta_{t+1}=\theta_t-\eta\,\bigl(\sign(c_t)+\lambda\theta_t\bigr)$}
\fgrow{norm}{$\ell_\infty$}
\fgrow{preconditioner}{none}
\fgrow{averages}{one EMA of $g$ (coefficient $\beta_2$), read out with weight $1-\beta_1$ on the current gradient}
\fgrow{state}{one number per parameter}
\fgrow{in the wild}{$\beta_1=0.9$, $\beta_2=0.99$, a learning rate several times smaller than AdamW's (\cref{ex:sign:lion})}
\fgrow{variants}{$\beta_1=\beta_2$ is Signum; with weight decay it solves a box-constrained problem \parencite{ChenLiuLiangLiu2024}}
\end{fieldguide}

\begin{fieldguide}{Adam}{\textcite{KingmaBa2015}}
\fgrow{update}{$\theta_{t+1}=\theta_t-\eta\,\hat m_t/(\sqrt{\hat v_t}+\epsilon)$; full entry in \cref{ch:adaptive}}
\fgrow{norm}{$\ell_\infty$ if $\beta_1=\beta_2=\epsilon=0$ (\cref{prop:sign:adam0}); otherwise steps of at most $C(\beta_1,\beta_2)\,\eta$ per coordinate (\cref{prop:sign:adambound})}
\fgrow{preconditioner}{diagonal, $\diag(\sqrt{\hat v_t})+\epsilon$}
\fgrow{averages}{EMA of $g$ ($\beta_1$) and of $g^2$ ($\beta_2$)}
\fgrow{state}{two numbers per parameter}
\fgrow{in the wild}{$\beta_1=0.9$, $\beta_2\in\{0.95,0.999\}$, $\epsilon=10^{-8}$}
\end{fieldguide}

\paragraph{What was thrown away.} The idealization keeps the step and discards the rest: the moving averages of Adam, Lion and Signum (they return only through the size of Adam's step and a heuristic for its mean, and in full in \cref{ch:momentum,ch:adaptive}), Adam's $\epsilon$, the variation of curvature along a trajectory, which a global bound in the coordinates replaces, and every feature of the noise except its variance. These are the right first simplifications because the sign is the one ingredient that all of these optimizers share and that gradient descent lacks.

\paragraph{What one had to notice.} Without memory, dividing by $\sqrt v$ is dividing by $\abs g$, so Adam's update is a vector of signs, a corner of the unit ball of the max norm. Once this is seen, Adam, Rprop, signSGD and Lion are one method in one geometry, and the questions are those of \cref{ch:steepest}: which dual norm measures the gradient, and which smoothness constant pays for the step. \Textcite{BernsteinNewhouse2024} put this observation at the center of a classification of optimizers by norms.

\section{The theory: the max-norm step}\label{sec:sign:theory}

\begin{proposition}[Adam without memory is sign descent]\label[proposition]{prop:sign:adam0}
Let $\epsilon\ge0$ and write Adam's update as $\theta_{t+1}=\theta_t-\eta u_t$, with $u_{t,i}=\hat m_{t,i}/(\sqrt{\hat v_{t,i}}+\epsilon)$ and $0/0=0$. (i) If $\beta_1=\beta_2=0$, then $u_{t,i}=\frac{\abs{g_{t,i}}}{\abs{g_{t,i}}+\epsilon}\sign(g_{t,i})$ for every $t$, with or without bias correction, and $u_t=\sign(g_t)$ if $\epsilon=0$; the same holds for RMSProp, $\theta_{t+1}=\theta_t-\eta g_t/(\sqrt{v_t}+\epsilon)$, with coefficient $0$. (ii) With bias correction, $u_0$ is given by the same formula for all $\beta_1,\beta_2\in[0,1)$: if $\epsilon=0$, Adam's first step is a sign step whatever the size of the first gradient.
\end{proposition}

\begin{proof}
(i) With $\beta_1=\beta_2=0$, $m_t=g_t$ and $v_t=g_t^2$, and the correction factors $1-0^{t+1}$ equal $1$, so $u_{t,i}=g_{t,i}/(\abs{g_{t,i}}+\epsilon)$; RMSProp with coefficient $0$ has the same update. (ii) Here $m_0=(1-\beta_1)g_0$ and $v_0=(1-\beta_2)g_0^2$, so $\hat m_0=g_0$ and $\hat v_0=g_0^2$.
\end{proof}

The factor $\abs g/(\abs g+\epsilon)$ makes memoryless Adam sign descent where the gradient is much larger than $\epsilon$ and gradient descent with step size $\eta/\epsilon$ where it is much smaller (\cref{ex:sign:epsilon}). With $\epsilon=0$, memoryless Adam, memoryless RMSProp and signSGD are one algorithm, and its guarantee is the descent inequality of \cref{ch:steepest} in the max norm.

\begin{theorem}[Sign descent in the max norm]\label{thm:sign:descent}
Let \cref{ass:sign:smooth} hold and let $\theta_{t+1}=\theta_t-\eta_t\sign(\nabla\Loss(\theta_t))$ with $\eta_t>0$. Then for every $t$
\begin{equation}\label{eq:sign:onestep}
  \Loss(\theta_{t+1})\le\Loss(\theta_t)-\eta_t\norm{\nabla\Loss(\theta_t)}_1+\tfrac12\eta_t^2\norm{\vec L}_1,
\end{equation}
where $\norm{\vec L}_1$ may be replaced by the sum of the $L_i$ with $\partial_i\Loss(\theta_t)\ne0$; and for every $T\ge1$
\begin{equation}\label{eq:sign:rate}
  \frac{\sum_{t<T}\eta_t\norm{\nabla\Loss(\theta_t)}_1}{\sum_{t<T}\eta_t}\le\frac{\Delta_0+\frac12\norm{\vec L}_1\sum_{t<T}\eta_t^2}{\sum_{t<T}\eta_t}.
\end{equation}
For a constant step size $\eta$ the right side is $\Delta_0/(\eta T)+\frac12\eta\norm{\vec L}_1$; if $\norm{\vec L}_1>0$ and $\eta=\sqrt{2\Delta_0/(\norm{\vec L}_1T)}$, it is $\sqrt{2\Delta_0\norm{\vec L}_1/T}$, which therefore bounds $\min_{t<T}\norm{\nabla\Loss(\theta_t)}_1$.
\end{theorem}

\begin{proof}
The plan is to apply \cref{ass:sign:smooth} to the step $u=-\eta_t\sign(\nabla\Loss(\theta_t))$ and to sum. Fix $t$, and write $g=\nabla\Loss(\theta_t)$ and $s=\sign(g)$. The linear term of the assumption is $\ip{g,-\eta_ts}=-\eta_t\sum_i\abs{g_i}=-\eta_t\norm{g}_1$. The quadratic term is $\frac12\eta_t^2\sum_iL_is_i^2$, and $s_i^2$ is $1$ where $g_i\ne0$ and $0$ elsewhere. This proves \eqref{eq:sign:onestep}. Rearranged, \eqref{eq:sign:onestep} bounds $\eta_t\norm{\nabla\Loss(\theta_t)}_1$ by $\Loss(\theta_t)-\Loss(\theta_{t+1})+\frac12\eta_t^2\norm{\vec L}_1$; summed over $t<T$, the loss differences telescope to $\Loss(\theta_0)-\Loss(\theta_T)\le\Delta_0$, and dividing by $\sum_t\eta_t$ gives \eqref{eq:sign:rate}. The stated constant $\eta$ minimizes $\Delta_0/(\eta T)+\frac12\eta\norm{\vec L}_1$, and a minimum is at most an average. If $\norm{\vec L}_1=0$, the assumption puts $\Loss$ below each of its tangent planes, so $\Loss$ is concave; being bounded below on $\R^d$, it is constant, and there is nothing to prove.
\end{proof}

The comparison with gradient descent comes from reading the same assumption in three geometries.

\begin{corollary}[Three guaranteed decreases]\label[corollary]{cor:sign:compare}
Let \cref{ass:sign:smooth} hold with every $L_i>0$, and let $g=\nabla\Loss(\theta)\ne0$. The gradient step $-g/L_{\max}$, the sign step $-(\norm{g}_1/\norm{\vec L}_1)\sign(g)$ and the diagonally preconditioned step $-\diag(\vec L)^{-1}g$ decrease $\Loss$ by at least $\norm{g}_2^2/(2L_{\max})$, $\norm{g}_1^2/(2\norm{\vec L}_1)$ and $\sum_ig_i^2/(2L_i)$ respectively. The third is at least the second, with equality if and only if $\abs{g_i}/L_i$ does not depend on $i$; the second is at least the first if and only if
\begin{equation}\label{eq:sign:density}
  \frac{\norm{g}_1^2}{\norm{g}_2^2}\ \ge\ \frac{\norm{\vec L}_1}{L_{\max}} .
\end{equation}
\end{corollary}

\begin{proof}
The plan is to evaluate \cref{ass:sign:smooth} at the three steps and to compare the results. The first and third bounds are the assumption at $u=-g/L_{\max}$, with $\sum_iL_iu_i^2\le L_{\max}\norm{u}_2^2$, and at $u_i=-g_i/L_i$. The second is \eqref{eq:sign:onestep} at $\eta=\norm{g}_1/\norm{\vec L}_1$, which maximizes $\eta\norm{g}_1-\frac12\eta^2\norm{\vec L}_1$. Writing $\abs{g_i}=(\abs{g_i}/\sqrt{L_i})\sqrt{L_i}$, the Cauchy--Schwarz inequality gives $\norm{g}_1^2\le\norm{\vec L}_1\sum_ig_i^2/L_i$, with equality exactly when $\abs{g_i}/\sqrt{L_i}$ is proportional to $\sqrt{L_i}$. The last claim is a rearrangement.
\end{proof}

Both sides of \eqref{eq:sign:density} lie between $1$ and $d$. The left side counts the coordinates over which the gradient is spread; the right side is an \term{effective rank}\index{effective rank} of the curvature in the coordinates of the parameters, which for a quadratic with diagonal Hessian and $\vec L=\diag(H)$ is $\tr H/\lmax$. This is the reading of the comparison of normalized steps in \cref{prop:steepest:whichnorm}, which has the exact constant $\norm{H}_{\infty\to1}$ in place of $\norm{\vec L}_1$: sign descent's guarantee is the better one when the gradient is spread over more coordinates than the curvature has stiff directions. A network has millions of parameters, most of which receive gradient, and the Hessians of trained networks have a few large outlier eigenvalues above a bulk of small ones \parencite{GhorbaniKrishnanXiao2019}; were such a Hessian diagonal in the parameters' coordinates, the advantage could approach $d$ over the effective rank. Whether it is close to diagonal is the subject of \cref{prop:sign:rotation}. The equality case says when the sign is as good as the best diagonal preconditioner: for a quadratic $\abs{g_i}/L_i$ is the length of the Newton step in coordinate $i$, so this happens when all coordinates are equally far from their targets.

Is the cost $\frac12\eta^2\norm{\vec L}_1$ an artifact of the proof? On a quadratic it is exact.

\begin{proposition}[The floor of sign descent on a quadratic]\label[proposition]{prop:sign:floor}
Let $\Loss(\theta)=\frac12(\theta-\theta^\star)^\top H(\theta-\theta^\star)$ with $H\succeq0$, run sign descent with a constant step size $\eta$, and write $g_t=\nabla\Loss(\theta_t)$ and $s_t=\sign(g_t)$.
\begin{enumerate}
\item For every $T\ge1$,
\begin{equation}\label{eq:sign:identity}
  \frac1T\sum_{t<T}\Bigl(\norm{g_t}_1-\frac\eta2\,s_t^\top Hs_t\Bigr)=\frac{\Delta_0-\Delta_T}{\eta T},
\end{equation}
and $\Delta_t$ is bounded; so $\frac1T\sum_{t<T}\norm{g_t}_1$ and $\frac1T\sum_{t<T}\frac12\eta s_t^\top Hs_t$ have the same limit points, all at most $\frac12\eta\norm{H}_{\infty\to1}$.
\item If $H=\diag(h_1,\dots,h_d)$ with all $h_i>0$, let $r_i\in[0,\eta)$ be the remainder of $\abs{\theta_{0,i}-\theta^\star_i}$ modulo $\eta$. From step $\lfloor\abs{\theta_{0,i}-\theta^\star_i}/\eta\rfloor$ on, the distance of coordinate $i$ from $\theta^\star_i$ alternates between $r_i$ and $\eta-r_i$ if $r_i>0$, and is $0$ if $r_i=0$. Hence $\frac1T\sum_{t<T}\norm{g_t}_1\to\frac12\eta\sum_{i:\,r_i>0}h_i$, which is $\frac12\eta\tr H$ unless some $\abs{\theta_{0,i}-\theta^\star_i}$ is a multiple of $\eta$.
\end{enumerate}
\end{proposition}

\begin{proof}
The plan is to use that the descent inequality is an identity for a quadratic, and to solve the separable case one coordinate at a time. (i) The identity is $\Loss(\theta_t-\eta s_t)=\Loss(\theta_t)-\eta\norm{g_t}_1+\frac12\eta^2s_t^\top Hs_t$; summing over $t<T$ and dividing by $\eta T$ gives \eqref{eq:sign:identity}. For boundedness, let $H\ne0$, let $\lambda_+$ be its smallest positive eigenvalue and $L_\infty=\norm{H}_{\infty\to1}$; then $s^\top Hs\le L_\infty$ for every $s\in\{-1,0,1\}^d$, since the convex function $x\mapsto x^\top Hx$ is maximized over the cube at a vertex. As $g_t$ lies in the range of $H$, $\Delta_t=\frac12g_t^\top H^+g_t\le\norm{g_t}_1^2/(2\lambda_+)$. If $\norm{g_t}_1\ge\frac12\eta L_\infty$ the identity gives $\Delta_{t+1}\le\Delta_t$, and otherwise $\Delta_{t+1}\le\Delta_t+\frac12\eta^2L_\infty<\eta^2L_\infty^2/(8\lambda_+)+\frac12\eta^2L_\infty$; by induction $\Delta_t$ never exceeds the larger of $\Delta_0$ and this number.
(ii) Coordinate $i$ evolves on its own: $x_t=\theta_{t,i}-\theta^\star_i$ satisfies $x_{t+1}=x_t-\eta\sign(x_t)$, so $\abs{x_{t+1}}=\bigl\lvert\abs{x_t}-\eta\bigr\rvert$ if $x_t\ne0$, and $x_{t+1}=0$ if $x_t=0$. The distance falls by $\eta$ per step for $\lfloor\abs{x_0}/\eta\rfloor$ steps and reaches $r_i$; if $r_i=0$ it stays there, and if $r_i>0$ the next distances are $\eta-r_i$ and $\abs{\eta-r_i-\eta}=r_i$. Over each period $h_i\abs{x_t}$ averages $\frac12\eta h_i$.
\end{proof}

For a separable quadratic the best constants are $\vec L=\diag(H)$, so the long-run average of $\norm{\nabla\Loss}_1$ equals the floor $\frac12\eta\norm{\vec L}_1$ of \cref{thm:sign:descent}, which cannot be improved. At a constant step size sign descent ends on a \term{two-cycle}\index{two-cycle}, the many-dimensional version of the prologue's example (\cref{ch:prologue}); only a decreasing schedule removes the floor (\cref{ch:schedules}).

Gradient descent does not depend on the orthonormal basis in which the parameters are written; sign descent does.

\begin{proposition}[Sign descent is not rotation invariant]\label[proposition]{prop:sign:rotation}
Let $H\succeq0$ have largest eigenvalue $\lmax$ with unit eigenvector $v_1$, and for an orthogonal $Q$ let $\Loss_Q(\phi)=\Loss(Q\phi)$, whose Hessian $Q^\top HQ$ has largest eigenvalue $\lmax$. Then (i) $\tr H\le\norm{Q^\top HQ}_{\infty\to1}\le d\lmax$, with the lower bound attained when the columns of $Q$ are eigenvectors of $H$ and the upper bound when $Q^\top v_1=\mathbf1/\sqrt d$; (ii) $\norm{Q^\top HQ}_{\infty\to1}\ge\lmax\norm{Q^\top v_1}_1^2$; (iii) if $Q$ is uniformly distributed on the orthogonal group, $\E\norm{Q^\top v_1}_1^2=1+\frac2\pi(d-1)$. So rotations change the $\ell_\infty$ smoothness constant by factors up to $d\lmax/\tr H$, and a random rotation multiplies it on average by at least $(1+\frac2\pi(d-1))\lmax/\tr H$.
\end{proposition}

\begin{proof}
The plan is to bound $s^\top As$, for $A=Q^\top HQ$, by averaging over random signs, by the top eigenvalue, and at the signs of the top eigenvector; recall $\norm{A}_{\infty\to1}=\max_{s\in\{\pm1\}^d}s^\top As$. (i) For independent uniform signs, $\E[s^\top As]=\sum_{i,j}A_{ij}\E[s_is_j]=\tr A=\tr H$, so some sign vector reaches $\tr H$; and $s^\top As\le\lmax\norm{s}_2^2=d\lmax$. If the columns of $Q$ are eigenvectors, $A$ is diagonal and nonnegative and $s^\top As=\tr A$ for all $s$. Some orthogonal $Q$ maps $\mathbf1/\sqrt d$ to $v_1$, and for it (ii) gives $\norm{A}_{\infty\to1}\ge d\lmax$. (ii) Let $s=\sign(Q^\top v_1)$, zeros replaced by $1$, and let $(\lambda_j,v_j)$ be the eigenpairs of $H$, all $\lambda_j\ge0$. Then $s^\top As=\sum_j\lambda_j\ip{v_j,Qs}^2\ge\lmax\ip{Q^\top v_1,s}^2=\lmax\norm{Q^\top v_1}_1^2$. (iii) The vector $u=Q^\top v_1$ is uniform on the sphere, so it has the law of $z/\norm z_2$ for $z\sim\cN(0,I_d)$ and is independent of $\norm z_2$. For $i\ne j$, $\frac2\pi=\E\abs{z_iz_j}=\E\norm z_2^2\,\E\abs{u_iu_j}=d\,\E\abs{u_iu_j}$; hence $\E\norm u_1^2=\sum_i\E u_i^2+\sum_{i\ne j}\E\abs{u_iu_j}=1+d(d-1)\frac{2}{\pi d}$.
\end{proof}

A random rotation also spreads the gradient, $\norm{Q^\top g}_1^2\approx\frac2\pi d\norm{g}_2^2$, so after it both sides of the criterion \eqref{eq:sign:density}, in its exact form, are about $\frac2\pi d$ or more on the right: the advantage is gone, not reversed. The guarantee of sign descent is a property of the coordinate system, and \cref{fig:sign:rotation} shows the dynamics following it.

\section{The theory: noise and memory}\label{sec:sign:noise}

A stochastic gradient has the wrong sign in some coordinates, and a sign method moves those coordinates the wrong way by a full step. The cost is a question about one random variable.

\begin{lemma}[Sign errors]\label[lemma]{lem:sign:signerror}
Let $X$ be a random variable with mean $\mu\ne0$ and variance at most $s^2$, where $s>0$. (i) $\Pr[\sign X\ne\sign\mu]\le s^2/(s^2+\mu^2)$. (ii) $\mu\,\E[\sign X]\ge\abs{\mu}-2\abs{\mu}s^2/(s^2+\mu^2)\ge\abs{\mu}-s$. (iii) If $X-\mu$ is symmetric about $0$ and unimodal, then $\mu\,\E[\sign X]\ge\min\bigl\{\mu^2/(\sqrt3\,s),\ \frac23\abs{\mu}\bigr\}>0$.
\end{lemma}

\begin{proof}
The plan is Cantelli's inequality for (i) and (ii) and Gauss's inequality for (iii). Replacing $X$ by $-X$ if necessary, let $\mu>0$ and write $\zeta=X-\mu$. (i) The event is $\{\zeta\le-\mu\}$, which for every $c\ge0$ lies in $\{(c-\zeta)^2\ge(c+\mu)^2\}$; by Markov's inequality its probability is at most $(s^2+c^2)/(c+\mu)^2$, which is $s^2/(s^2+\mu^2)$ at $c=s^2/\mu$ (Cantelli's inequality\index{Cantelli's inequality}). (ii) Pointwise $\sign X\ge1-2\,\ind[X\le0]$; take expectations and use (i). The second inequality is $2\mu s\le s^2+\mu^2$. (iii) A unimodal law with mode $0$ has no atoms away from $0$, and by symmetry $\Pr[\zeta<-\mu]=\Pr[\zeta>\mu]$, so $\E[\sign X]=\Pr[\zeta>-\mu]-\Pr[\zeta>\mu]=\Pr[\abs\zeta<\mu]$. Gauss's inequality\index{Gauss's inequality} (\cref{lem:toolkit:gauss}), for a unimodal law with mode $0$ and second moment $\tau^2\le s^2$, bounds $\Pr[\abs\zeta\ge k]$ by $1-k/(\sqrt3\tau)$ if $k\le2\tau/\sqrt3$ and by $4\tau^2/(9k^2)$ otherwise. Both bounds increase with $\tau$, so they hold with $s$ for $\tau$, and at $k=\mu$ they give $\Pr[\abs\zeta<\mu]\ge\mu/(\sqrt3s)$ if $\mu\le2s/\sqrt3$ and $\Pr[\abs\zeta<\mu]\ge1-\frac{4s^2}{9\mu^2}\ge\frac23$ otherwise.
\end{proof}

By (ii), the expected progress $\mu\E\sign X$ falls short of the noiseless $\abs\mu$ by at most the noise level $s$. Summed over coordinates, this is the stochastic version of \cref{thm:sign:descent}.

\begin{theorem}[signSGD]\label{thm:sign:signsgd}
Let \cref{ass:sign:smooth,ass:sign:noise} hold, and run signSGD, $\theta_{t+1}=\theta_t-\eta_t\sign(g_t)$, with deterministic $\eta_t>0$. Then for every $T\ge1$
\[
  \frac{\sum_{t<T}\eta_t\,\E\norm{\nabla\Loss(\theta_t)}_1}{\sum_{t<T}\eta_t}\le\frac{\Delta_0+\frac12\norm{\vec L}_1\sum_{t<T}\eta_t^2}{\sum_{t<T}\eta_t}+\frac{\norm{\vec\sigma}_1}{\sqrt B}.
\]
For a constant step size the right side is $\Delta_0/(\eta T)+\frac12\eta\norm{\vec L}_1+\norm{\vec\sigma}_1/\sqrt B$. With $\eta=\sqrt{2\Delta_0/(\norm{\vec L}_1T)}$ and $B=T$, after $N=BT$ stochastic gradients, it is $\bigl(\sqrt{2\Delta_0\norm{\vec L}_1}+\norm{\vec\sigma}_1\bigr)N^{-1/4}$.
\end{theorem}

\begin{proof}
The plan is the proof of \cref{thm:sign:descent} in conditional expectation, with \cref{lem:sign:signerror}(ii) for the linear term. Let $\cF_t$ be generated by $\theta_0,g_0,\dots,g_{t-1}$, and write $\mu=\nabla\Loss(\theta_t)$. \Cref{ass:sign:smooth} at $u=-\eta_t\sign(g_t)$ gives
\[
  \E[\Loss(\theta_{t+1})\mid\cF_t]\le\Loss(\theta_t)-\eta_t\sum_{i=1}^d\mu_i\,\E[\sign g_{t,i}\mid\cF_t]+\tfrac12\eta_t^2\norm{\vec L}_1 .
\]
Conditionally on $\cF_t$, the coordinate $g_{t,i}$ has mean $\mu_i$ and variance at most $\sigma_i^2/B$, so \cref{lem:sign:signerror}(ii) gives $\mu_i\E[\sign g_{t,i}\mid\cF_t]\ge\abs{\mu_i}-\sigma_i/\sqrt B$ if $\mu_i\ne0$ and $\sigma_i>0$; the bound also holds if $\mu_i=0$, when the left side vanishes, and if $\sigma_i=0$, when $g_{t,i}=\mu_i$. Sum over $i$, take expectations, sum over $t<T$, use $\E\Loss(\theta_T)\ge\Loss^\star$ and divide by $\sum_t\eta_t$. With the stated $\eta$ and $B=T$ the first two terms equal $\sqrt{2\Delta_0\norm{\vec L}_1/T}$, the third is $\norm{\vec\sigma}_1/\sqrt T$, and $\sqrt T=N^{1/4}$.
\end{proof}

With a batch that grows with the horizon, signSGD reaches the rate $N^{-1/4}$ of stochastic gradient descent, with $\ell_1$ norms of the gradient and of the noise in place of $\ell_2$ norms; this is the comparison of \textcite{BernsteinEtAl2018}, whose noise term has the factor $2$ of Markov's inequality where Cantelli's gives $1$. At a fixed batch size, however, the term $\norm{\vec\sigma}_1/\sqrt B$ survives every step size. This is not a weakness of the proof: the sign of a minibatch gradient is right in most batches when the median of the gradient has the sign of its mean, and a rare event can separate the two.

\begin{proposition}[A rare class defeats signSGD]\label[proposition]{prop:sign:rareclass}
Let labels $y\in\{0,1\}$ be independent with $\Pr[y=1]=\pi\in(0,\frac12)$, let a model predict $\varsigma(b)=1/(1+e^{-b})$ with one bias $b\in\R$, and let $\Loss(b)=-\pi\log\varsigma(b)-(1-\pi)\log(1-\varsigma(b))$. SignSGD on batches of $B$ labels is $b_{t+1}=b_t-\eta\,\sign(\varsigma(b_t)-\bar y_t)$, where $\bar y_t$ is the fraction of ones in batch $t$. (i) $\Loss$ satisfies \cref{ass:sign:smooth,ass:sign:noise} with $L_1=\frac14$ and $\sigma_1^2=\pi(1-\pi)$, and is minimized where $\varsigma(b^\star)=\pi$. (ii) If $q=(1-\pi)^B>\frac12$, then for every $\eta>0$ and $b_0$, almost surely $\limsup_tb_t/t\le-\eta(2q-1)<0$; hence $\Loss(b_t)\to\infty$ and $\nabla\Loss(b_t)\to-\pi$.
\end{proposition}

\begin{proof}
(i) Here $\Loss'(b)=\varsigma(b)-\pi$ and $0<\Loss''(b)=\varsigma(b)(1-\varsigma(b))\le\frac14$, and the minibatch gradient $\varsigma(b)-\bar y_t$ averages $B$ independent terms $\varsigma(b)-y$ of mean $\varsigma(b)-\pi$ and variance $\pi(1-\pi)$.
(ii) The plan is to bound the drift of $b_t$ and to control its fluctuations by a martingale argument. With probability $q$ batch $t$ has no label $1$; then $\bar y_t=0<\varsigma(b_t)$ and the step is $-\eta$. A step $+\eta$ needs $\bar y_t>\varsigma(b_t)$, which has probability at most $1-q$. So $D_t=\E[b_{t+1}-b_t\mid b_0,\dots,b_t]\le-\eta(2q-1)$ whatever $b_t$ is. The differences $M_t=b_{t+1}-b_t-D_t$ are martingale differences with $\abs{M_t}\le2\eta$, and by the Azuma--Hoeffding inequality (\cref{thm:toolkit:azuma}) $\Pr\bigl[\abs{\sum_{s<t}M_s}\ge\eps t\bigr]\le2e^{-\eps^2t/(8\eta^2)}$ for each $\eps>0$. These are summable in $t$, so by the Borel--Cantelli lemma $\frac1t\sum_{s<t}M_s\to0$ almost surely, and $b_t/t=b_0/t+\frac1t\sum_{s<t}(D_s+M_s)$ has $\limsup$ at most $-\eta(2q-1)$. Then $\varsigma(b_t)\to0$, $\Loss(b_t)\ge-\pi\log\varsigma(b_t)\to\infty$ and $\Loss'(b_t)\to-\pi$.
\end{proof}

The batches without the rare label are the majority, and in them the gradient is positive, so the sign follows the majority. The mean gradient $\varsigma(b)-\pi$ is negative wherever the predicted probability is too small, but the median of the minibatch gradient is $\varsigma(b)>0$ for every $b$: signSGD descends the median, which here has no zero.\sidenote{In one dimension this is classical: with $g=\theta-\xi$, signSGD is stochastic subgradient descent on $\E\abs{\theta-\xi}$, whose minimizers are the medians of $\xi$.} Once $\varsigma(b_t)<1/B$, every batch with a label $1$ pushes $b$ up, and the drift is exactly $-\eta(2q-1)$ (\cref{fig:sign:noise}). The condition $(1-\pi)^B>\frac12$ holds for every class rarer than $1-2^{-1/B}\approx\ln2/B$, and a softmax layer over heavy-tailed classes, such as tokens, has many (\cref{ch:transformers}). \Textcite{ReddiKaleKumar2018} built their counterexample to Adam with a small $\beta_2$ from the same ingredient, a large gradient that comes rarely against small ones of the opposite sign that come often (\cref{ch:adaptive}). The proposition also shows that the constant in \cref{thm:sign:signsgd} cannot be lowered: as $\eta\to0$ and $T\to\infty$ the average of $\abs{\nabla\Loss(b_t)}$ is $\pi$ and the bound is $\sqrt{\pi(1-\pi)/B}$, and their ratio tends to $1$ for $B=1$ and $\pi\uparrow\frac12$, and to $\sqrt{\ln2}\approx0.83$ for large $B$ and $\pi$ just below $1-2^{-1/B}$. For symmetric noise the median is the mean, and the bias disappears.

\begin{proposition}[Symmetric unimodal noise: no bias]\label[proposition]{prop:sign:symmetric}
Let \cref{ass:sign:smooth,ass:sign:noise} hold with every $\sigma_i>0$, and let each $\zeta_{t,i}$ be, conditionally on the past, symmetric about $0$ and unimodal. With $\Phi_B(\mu)=\sum_i\min\{\sqrt B\mu_i^2/(\sqrt3\sigma_i),\frac23\abs{\mu_i}\}$, signSGD with a constant step size $\eta$ satisfies
\begin{equation}\label{eq:sign:mixed}
  \frac1T\sum_{t<T}\E\,\Phi_B\bigl(\nabla\Loss(\theta_t)\bigr)\le\frac{\Delta_0}{\eta T}+\frac\eta2\norm{\vec L}_1 .
\end{equation}
\end{proposition}

\begin{proof}
Repeat the proof of \cref{thm:sign:signsgd} with \cref{lem:sign:signerror}(iii), for $s=\sigma_i/\sqrt B$, in place of (ii): the one-step inequality becomes $\E[\Loss(\theta_{t+1})\mid\cF_t]\le\Loss(\theta_t)-\eta\,\Phi_B(\nabla\Loss(\theta_t))+\frac12\eta^2\norm{\vec L}_1$, both sides of the per-coordinate bound being $0$ when $\mu_i=0$, and the sum telescopes.
\end{proof}

There is no bias term: with $\eta\propto T^{-1/2}$ the right side of \eqref{eq:sign:mixed} is $O(T^{-1/2})$ at every batch size. The batch size enters only the measure of progress, which is $\frac23\abs{\mu_i}$ in coordinates whose gradient exceeds the noise $\sigma_i/\sqrt B$ by a constant factor, and $\sqrt B\mu_i^2/(\sqrt3\sigma_i)$ in coordinates dominated by noise. In the second regime signSGD is a preconditioned gradient method.

\begin{heuristic}[Heuristic derivation of the effective step]
Let the noise be Gaussian and $\abs{\mu_i}\ll\sigma_i/\sqrt B$. Then $\Var(\sign g_{t,i})\approx1$ and
\[
  \E[\sign g_{t,i}]=\erf\Bigl(\frac{\sqrt B\,\mu_i}{\sqrt2\,\sigma_i}\Bigr)\approx\sqrt{\frac{2B}\pi}\,\frac{\mu_i}{\sigma_i},
\]
so a step of signSGD has the mean and twice the variance, up to the factor $\pi/2$ in place of $2$, of a step of stochastic gradient descent with the step size $\alpha_i=\eta\sqrt{2B/\pi}/\sigma_i$: signSGD is gradient descent preconditioned by $\diag(\sigma_i)^{-1}$, with noise inflated by $\sqrt{\pi/2}$. On $\Loss(\theta)=\frac12h\theta^2$, the stationary variance $\alpha\frac\pi2\frac{\sigma^2}{B}/(2h)$ of that method gives the floor
\begin{equation}\label{eq:sign:gaussfloor}
  \E\abs{\nabla\Loss}\approx(2\pi)^{-1/4}\sqrt{\eta h\sigma}\;B^{-1/4},
\end{equation}
which vanishes as $\eta\to0$. Keeping $\alpha_i\propto B$, the linear scaling rule of gradient descent (\cref{ch:noise}), requires $\eta\propto\sqrt B$: the \term{square-root scaling rule}\index{square-root scaling rule}, which \textcite{MalladiEtAl2022} derived for adaptive methods from stochastic differential equations.
\end{heuristic}

Memory changes the picture in two ways. Momentum averages gradients over time before the sign is taken, which moves the median towards the mean; \cref{ch:momentum} proves that normalized steepest descent with momentum converges in every norm, and \cref{fig:sign:noise} shows Signum learning the rare class that defeats signSGD. Adam's second moment changes the length of the step, and the first question is how long the step can get.

\begin{proposition}[The size of Adam's step]\label[proposition]{prop:sign:adambound}
Let $0\le\beta_1<1$, $0<\beta_2<1$, $\rho=\beta_1^2/\beta_2$ and $\epsilon=0$, and consider one coordinate of Adam without bias correction. For every $t\ge0$ and every gradient sequence with $v_t>0$,
\begin{equation}\label{eq:sign:adambound}
  \frac{\abs{m_t}}{\sqrt{v_t}}\le\frac{1-\beta_1}{\sqrt{1-\beta_2}}\Bigl(\sum_{k=0}^t\rho^k\Bigr)^{1/2},
\end{equation}
with equality if and only if $g_{t-k}=c\,(\beta_1/\beta_2)^k$ for $k=0,\dots,t$ and some $c\ne0$ (with $0^0=1$). Consequently, if $\beta_1^2<\beta_2$, the supremum over all $t$ and all gradient sequences is
\begin{equation}\label{eq:sign:C}
  C(\beta_1,\beta_2)=\frac{1-\beta_1}{\sqrt{(1-\beta_2)(1-\beta_1^2/\beta_2)}},
\end{equation}
and if $\beta_1^2\ge\beta_2$ the ratio is unbounded. With bias correction the bound \eqref{eq:sign:adambound} is multiplied by $\sqrt{1-\beta_2^{t+1}}/(1-\beta_1^{t+1})$; it is then $1$ at $t=0$ and below $C(\beta_1,\beta_2)$ for all $t$ (\cref{ex:sign:biascorr}).
\end{proposition}

\begin{proof}
The plan is to write $m_t$ as a weighted sum of gradients and to apply the Cauchy--Schwarz inequality with the weights of $v_t$. From $m_{-1}=v_{-1}=0$, $m_t=(1-\beta_1)\sum_{k=0}^t\beta_1^kg_{t-k}$ and $v_t=(1-\beta_2)\sum_{k=0}^t\beta_2^kg_{t-k}^2$. Splitting $\beta_1^kg_{t-k}=(\beta_1^k\beta_2^{-k/2})(\beta_2^{k/2}g_{t-k})$,
\[
  \Bigl\lvert\sum_{k=0}^t\beta_1^kg_{t-k}\Bigr\rvert\le\Bigl(\sum_{k=0}^t\rho^k\Bigr)^{1/2}\Bigl(\sum_{k=0}^t\beta_2^kg_{t-k}^2\Bigr)^{1/2}=\Bigl(\sum_{k=0}^t\rho^k\Bigr)^{1/2}\Bigl(\frac{v_t}{1-\beta_2}\Bigr)^{1/2},
\]
and multiplying by $1-\beta_1$ gives \eqref{eq:sign:adambound}. Equality holds exactly when $(\beta_2^{k/2}g_{t-k})_k$ is parallel to $(\beta_1^k\beta_2^{-k/2})_k$, that is, $g_{t-k}=c(\beta_1/\beta_2)^k$, with $c\ne0$ because $v_t>0$. If $\rho<1$ the sums $\sum_{k\le t}\rho^k$ increase to $1/(1-\rho)$, and the equality sequences make \eqref{eq:sign:C} the supremum; if $\rho\ge1$ the sums are unbounded. At $t=0$ the bias correction factor is $\sqrt{1-\beta_2}/(1-\beta_1)$, the reciprocal of the bound.
\end{proof}

The bound is $C=7.27$ at the defaults $(0.9,0.999)$, $2.35$ at $(0.9,0.99)$, $1.16$ at $(0.9,0.95)$, and $1$ whenever $\beta_1=\beta_2$. A single large gradient after a run of zeros gives only $(1-\beta_1)/\sqrt{1-\beta_2}$, which is $3.16$ at the defaults;\sidenote{\Textcite[Section~2.1]{KingmaBa2015} describe the effective step as approximately bounded by this single-spike value or by the learning rate. \Cref{prop:sign:adambound} gives the exact supremum.} the worst case is a gradient that grows by the factor $\beta_2/\beta_1$ at every step, so that the slow second moment lags behind the fast first moment. Adam is thus sign descent of its momentum with a bounded multiplier: each coordinate moves by $\eta$ times a number in $[0,C]$ in the direction of $-\sign(m_t)$, and $\norm{\theta_{t+1}-\theta_t}_\infty\le C\eta$, a trust region in the max norm. A smaller $\beta_2$ keeps the multiplier near $1$, and the condition $\beta_1<\sqrt{\beta_2}$ for a bounded step reappears in convergence analyses of Adam and AMSGrad \parencite{ReddiKaleKumar2018,ZhangEtAl2022}. The typical multiplier, as opposed to the worst one, is set by the signal-to-noise ratio.

\begin{heuristic}[Heuristic derivation of Adam's mean step]
Let the gradients of coordinate $i$ be independent over steps with mean $\mu_i$ and variance $s_i^2$, and let $\beta_2$ be so close to $1$ that $\hat v_{t,i}$, an average of about $1/(1-\beta_2)$ squared gradients, concentrates near $\mu_i^2+s_i^2$; $\hat m_{t,i}$ has mean $\mu_i$ and variance $(1-\beta_1)^2\sum_k\beta_1^{2k}s_i^2=\frac{1-\beta_1}{1+\beta_1}s_i^2$. Treating $\hat v_{t,i}$ as a constant,
\begin{equation}\label{eq:sign:snr}
  \E\Bigl[\frac{\hat m_{t,i}}{\sqrt{\hat v_{t,i}}}\Bigr]\approx\frac{\sign(\mu_i)}{\sqrt{1+s_i^2/\mu_i^2}},
\end{equation}
and the mean square of the update is about $\bigl(\mu_i^2+\frac{1-\beta_1}{1+\beta_1}s_i^2\bigr)/(\mu_i^2+s_i^2)$. The mean update is the sign shrunk by the factor $(1+s_i^2/\mu_i^2)^{-1/2}$ of \textcite{BallesHennig2018}: near $\sign(\mu_i)$ for a consistent gradient and near $\mu_i/s_i$ for noise, where Adam, like signSGD, is gradient descent preconditioned by the noise. The root mean square of the update lies between $\sqrt{(1-\beta_1)/(1+\beta_1)}=0.23$, for pure noise at $\beta_1=0.9$, and $1$.
\end{heuristic}

For $(\beta_1,\beta_2)=(0.9,0.999)$ and Gaussian gradients, simulation confirms \eqref{eq:sign:snr} and the mean square to within $0.002$ for signal-to-noise ratios from $0$ to $3$, and in $2\times10^6$ steps of pure noise the multiplier never exceeded $1.5$, far from the worst case $7.27$. SignSGD's mean update under Gaussian noise is exactly $\erf\bigl(\mu_i/(\sqrt2s_i)\bigr)$, within $0.07$ of \eqref{eq:sign:snr}: in this regime Adam is very nearly a sign method with a different nonlinearity.

\section{Reading the theory}\label{sec:sign:reading}

\paragraph{The guarantee is exact, and it belongs to the basis.} \Cref{fig:sign:rotation} runs gradient descent and sign descent on $\Loss(\theta)=\frac12\sum_ih_i\theta_i^2$ with $d=256$ and $h_i=i^{-2}$, so that $\tr H=1.64$ and the effective rank is below $2$, and on the same quadratic in randomly rotated coordinates. In panel (b) the long-run average of $\norm{\nabla\Loss}_1$ on the diagonal problem lies on the line $\frac12\eta\tr H$ of \cref{prop:sign:floor} to all printed digits. On the rotated problem it lies on $\frac12\eta\hat L_\infty$, where $\hat L_\infty=161.2$ is the largest $s^\top Q^\top HQs$ that a local search over sign vectors finds; the certified lower bound of \cref{prop:sign:rotation}(ii) is $160.9$, and its mean over rotations $163.3$. In every run the final two-cycle used the best sign vector of the search, so the floor $\frac12\eta\norm{\cdot}_{\infty\to1}$ of \cref{thm:sign:descent} is attained in both bases, and the rotation multiplies it by $98$, against the $\frac2\pi d\lmax/\tr H=99$ of \cref{prop:sign:rotation}.

\begin{figure}[t]
\begin{fullwidth}
\centering
\begin{tikzpicture}
\begin{groupplot}[group style={group size=2 by 1,horizontal sep=17mm},wplot,width=82mm,height=54mm]
\nextgroupplot[xmode=log,ymode=log,xlabel={step $t$},ylabel={loss $\Loss(\theta_t)-\Loss^\star$},xmin=1,xmax=1e4,ymin=1e-5,ymax=1,
  legend style={at={(0.02,0.02)},anchor=south west,font=\scriptsize},legend cell align=left]
\addplot[theory] table[x=t,y=gdexact] {figures/data/02-sign-rotation-loss.dat};
\addplot[theory3] table[x=t,y=signexact] {figures/data/02-sign-rotation-loss.dat};
\addplot[simulation2] table[x=t,y=gdrot] {figures/data/02-sign-rotation-loss.dat};
\addplot[simulation] table[x=t,y=signdiag] {figures/data/02-sign-rotation-loss.dat};
\addplot[only marks,mark=triangle*,mark size=2pt,color=sea,mark options={fill=sea,draw=white,line width=0.2pt}] table[x=t,y=signrot] {figures/data/02-sign-rotation-loss.dat};
\legend{gradient descent (exact; both bases),sign descent (exact; diagonal),gradient descent (rotated),sign descent (diagonal),sign descent (rotated)}
\nextgroupplot[xmode=log,ymode=log,xlabel={step size $\eta$},ylabel={long-run average of $\norm{\nabla\Loss(\theta_t)}_1$},xmin=7e-4,xmax=0.14,ymin=3e-4,ymax=20,
  legend style={at={(0.02,0.98)},anchor=north west,font=\scriptsize},legend cell align=left]
\addplot[theory] table[x=eta,y=diagtheory] {figures/data/02-sign-rotation-floor.dat};
\addplot[theory3] table[x=eta,y=rottheory] {figures/data/02-sign-rotation-floor.dat};
\addplot[simulation] table[x=eta,y=diagsim] {figures/data/02-sign-rotation-floor.dat};
\addplot[simulation2] table[x=eta,y=rotsim] {figures/data/02-sign-rotation-floor.dat};
\legend{$\frac12\eta\tr H$ (diagonal),$\frac12\eta\hat L_\infty$ (rotated),diagonal (run),rotated (run)}
\end{groupplot}
\end{tikzpicture}
\end{fullwidth}
\caption{Gradient descent and sign descent on a quadratic with curvatures $h_i=i^{-2}$, $d=256$, $\theta^\star=0$ and $\theta_0\sim\cN(0,I)$, in the eigenbasis and after one uniformly random rotation (seed in \texttt{02-sign-rotation.py}). Left: the loss of gradient descent with step size $1/\lmax$ and of sign descent with $\eta=0.01$; curves are exact closed forms, points are runs. Gradient descent is the same in both bases; sign descent is far ahead on the diagonal problem and is caught on a two-cycle with loss $2.1\times10^{-2}$ after about $300$ steps on the rotated one. Right: the average of $\norm{\nabla\Loss}_1$ over the final cycle against $\eta$, with the floor of \cref{prop:sign:floor} and the floor $\frac12\eta\hat L_\infty$ of \cref{thm:sign:descent} for the rotated Hessian, $\hat L_\infty$ found by local search.}
\label{fig:sign:rotation}
\end{figure}

\paragraph{A time scale without curvature.} On the diagonal problem coordinate $i$ reaches its two-cycle after $\lfloor\abs{\theta_{0,i}}/\eta\rfloor$ steps whatever its curvature, so sign descent with $\eta=0.01$ is at its floor, near $2\times10^{-5}$, after about $300$ steps. Gradient descent with step size $1/\lmax$ needs about $\lmax/h_i=i^2$ steps for coordinate $i$, and its expected loss $\frac12\sum_ii^{-2}(1-i^{-2})^{2t}$ decays like $\frac{\sqrt\pi}{4\sqrt{2t}}\approx0.31\,t^{-1/2}$ while $t\ll d^2$: it is $1.5\times10^{-2}$ at $t=300$ and $1.3\times10^{-3}$ at $t=10^4$. This is the regime in which \cref{cor:sign:compare} favors the sign: curvature of effective rank below $2$, a spread gradient, and distances of the same order in all coordinates.

\paragraph{The trap after a rotation.} On the rotated problem gradient descent is unchanged to rounding error. Sign descent makes fast progress for a hundred steps and is then caught on a two-cycle with loss $2.1\times10^{-2}$, a thousand times its diagonal floor, after which gradient descent is ahead. The centers of the two-cycles form a box in gradient space, which the inverse Hessian stretches along the directions of small curvature; when these are not coordinate axes, the box contains points of high loss (\cref{ex:sign:cycles}). A decreasing step size breaks the cycle; Adam, whose effective step changes over time, escapes it but still loses orders of magnitude to the rotation (\cref{ex:sign:adamrot}).

\paragraph{Median against mean.} \Cref{fig:sign:noise}(a) shows one coordinate of signSGD with $B=16$ under two noises of the same variance. The iterates move on a lattice of spacing $\eta$ and form a birth--death chain, whose stationary law is computed exactly. With Gaussian noise the floor follows \eqref{eq:sign:gaussfloor} and vanishes like $\sqrt\eta$, as \cref{prop:sign:symmetric} requires. With rare-event noise, in which a fraction $p=0.04$ of the examples push hard one way and the rest gently the other way, $(1-p)^B=0.52>\frac12$, and as $\eta\to0$ the floor rises to $c=\sigma\sqrt{p/(1-p)}=0.82\,\sigma/\sqrt B$, the gradient at the point where the median of the minibatch gradient vanishes, $82\%$ of the bias term of \cref{thm:sign:signsgd}. A smaller step makes it worse; the lowest floor is near $\eta=0.03$. Panel (b) is \cref{prop:sign:rareclass} for a class of frequency $\pi=0.01$: the bias drifts at exactly the predicted speed for $B=16,32,64$, the speed vanishes at the threshold $B=\ln2/(-\ln(1-\pi))=69$, and Signum with $\beta=0.9$ holds the bias at its optimum.

\begin{figure}[t]
\begin{fullwidth}
\centering
\begin{tikzpicture}
\begin{groupplot}[group style={group size=2 by 1,horizontal sep=17mm},wplot,width=82mm,height=54mm]
\nextgroupplot[xmode=log,ymode=log,xlabel={step size $\eta$},ylabel={stationary $\E\abs{\nabla\Loss(\theta_t)}$},xmin=8e-4,xmax=0.4,ymin=7e-3,ymax=0.6,
  legend style={at={(0.98,0.02)},anchor=south east,font=\scriptsize},legend cell align=left]
\addplot[guide] table[x=eta,y=bound] {figures/data/02-sign-noise-eta.dat};
\addplot[guide,sandline!55!black] table[x=eta,y=bias] {figures/data/02-sign-noise-eta.dat};
\addplot[theory] table[x=eta,y=gaussexact] {figures/data/02-sign-noise-eta.dat};
\addplot[theory2] table[x=eta,y=heuristic] {figures/data/02-sign-noise-eta.dat};
\addplot[theory3] table[x=eta,y=rareexact] {figures/data/02-sign-noise-eta.dat};
\addplot[simulation2] table[x=eta,y=gausssim] {figures/data/02-sign-noise-eta.dat};
\addplot[simulation] table[x=eta,y=raresim] {figures/data/02-sign-noise-eta.dat};
\legend{bound of \cref{thm:sign:signsgd},median gap $c$,Gaussian (exact),heuristic \eqref{eq:sign:gaussfloor},rare event (exact),Gaussian (runs),rare event (runs)}
\nextgroupplot[xlabel={step $t$},ylabel={bias $b_t$},xmin=0,xmax=20000,ymin=-150,ymax=8,xtick={0,5000,10000,15000,20000},
  xticklabels={$0$,$5000$,$10000$,$15000$,$20000$},legend style={at={(0.02,0.02)},anchor=south west,font=\scriptsize},legend cell align=left]
\addplot[theory] table[x=t,y=theory16] {figures/data/02-sign-noise-rare.dat};
\addplot[theory2] table[x=t,y=theory32] {figures/data/02-sign-noise-rare.dat};
\addplot[theory3] table[x=t,y=theory64] {figures/data/02-sign-noise-rare.dat};
\addplot[simulation,mark size=1.2pt] table[x=t,y=sim16] {figures/data/02-sign-noise-rare.dat};
\addplot[simulation,mark size=1.2pt,forget plot] table[x=t,y=sim32] {figures/data/02-sign-noise-rare.dat};
\addplot[simulation,mark size=1.2pt,forget plot] table[x=t,y=sim64] {figures/data/02-sign-noise-rare.dat};
\addplot[simulation2,mark size=1.3pt] table[x=t,y=signum16] {figures/data/02-sign-noise-rare.dat};
\legend{drift for $B=16$,drift for $B=32$,drift for $B=64$,signSGD (runs),Signum{,} $B=16$ (runs)}
\end{groupplot}
\end{tikzpicture}
\end{fullwidth}
\caption{SignSGD and the shape of the noise. Left: one coordinate, $\Loss(\theta)=\frac12h\theta^2$ with $h=\sigma=1$ and $B=16$, under Gaussian per-example noise and under the rare-event noise $\sigma(X-p)/\sqrt{p(1-p)}$ with $X$ Bernoulli($p$), $p=0.04$. Curves: exact stationary values of the lattice chain; points: $1000$ runs averaged over at least $4000$ steps after a burn-in; dashed: the bound $\frac12\eta h+\sigma/\sqrt B$ and the median gap $c$. Right: the bias of a class of frequency $\pi=0.01$ under signSGD with $\eta=0.01$ from $b^\star$, averaged over $200$ runs; lines are the exact drift $-\eta(2(1-\pi)^B-1)t$; squares are Signum ($\beta=0.9$). Script: \texttt{02-sign-noise.py}.}
\label{fig:sign:noise}
\end{figure}

\section{What the theory misses}\label{sec:sign:misses}

\paragraph{What was idealized away.} Memory first: the theorems are about the memoryless step, Adam's averages enter only through a bound on its step and a heuristic for its mean, and nothing here proves that Adam converges or explains why $\beta_1=0.9$ and $\beta_2$ near $1$ work (\cref{ch:momentum,ch:adaptive}). Adam's $\epsilon$ is zero except in \cref{ex:sign:epsilon}. The curvature is bounded by one vector $\vec L$ everywhere, a strong assumption for networks, whose sharpness adapts to the step size (\cref{ch:eos}); the proofs use $\vec L$ only along the trajectory, but nothing tells us its value in advance. The noise is described by its variance, yet the outcome depends on its shape: the same variance gives no bias when the noise is symmetric and divergence when it is skewed. And most statements use a constant step size, while practice decays it (\cref{ch:schedules}).

\paragraph{Where it breaks.} (1) SignSGD without momentum fails at small batch sizes: the rare class of \cref{prop:sign:rareclass} is convex and smooth, and its loss diverges. \Textcite{KarimireddyEtAl2019} give further examples in which signSGD does not converge, and \textcite{SafaryanRichtarik2021} prove convergence under a condition on the probability that each sign is correct, which fails exactly when the median and the mean of a coordinate's gradient have different signs. (2) Adam with long memories is not sign descent: at $\beta_2=0.999$ its multiplier can reach $7.27$, and typically it is a smooth function of the signal-to-noise ratio \eqref{eq:sign:snr}, so the theorems for sign descent do not transfer to it. (3) \Cref{cor:sign:compare} compares upper bounds, and worst-case rates do not explain the size of the gap between Adam and gradient descent on transformers, which \textcite{KunstnerEtAl2023} measure in full batch (\cref{ch:transformers}). (4) The advantage of the max norm depends on the basis (\cref{prop:sign:rotation}), and the theory does not say why the coordinates of trained networks are favorable. \Textcite{XieMohamadiLi2025} relate Adam's sensitivity to rotations to a smoothness constant of the $\ell_\infty$ type, and \textcite{ZhangEtAl2024} find transformer Hessians close to block diagonal, with blocks that correspond to groups of parameters; both are measurements, not derivations. (5) The floor of sign descent can be far worse than \cref{prop:sign:floor} suggests: after a rotation, the trapping two-cycle has a loss a thousand times higher (\cref{fig:sign:rotation}).

\paragraph{Repairs and open problems.} Momentum is the main repair: averaging before taking the sign reduces the noise and moves the median towards the mean, normalized steepest descent with momentum converges in every norm (\cref{ch:momentum}), and Signum and Lion train language models about as well as Adam \parencite{ZhaoEtAl2025}. Error feedback, which adds the part of the gradient lost to compression to the next gradient, makes sign compression converge \parencite{SeideEtAl2014,KarimireddyEtAl2019}; majority vote reduces the noise of the sign when it is symmetric and unimodal \parencite{BernsteinEtAl2019} (\cref{ex:sign:vote}). Analyses of Adam in the $\ell_\infty$ geometry \parencite{XieMohamadiLi2025}, its implicit bias towards $\ell_\infty$ max-margin solutions \parencite{ZhangZouCao2024} (\cref{ch:bias}), and the $\ell_\infty$-constrained problems that AdamW and Lion solve with weight decay \parencite{XieLi2024,ChenLiuLiangLiu2024} (\cref{ch:decay}) carry the max-norm picture to the full algorithms; Shampoo and SOAP remove the dependence on the basis by rotating into an estimated eigenbasis (\cref{ch:kronecker}). Open: why the coordinates of trained networks are good ones for sign methods (\cref{ex:sign:basis}); a convergence theory of Adam that uses both its geometry and its memory and explains its speed; and the batch sizes and momenta at which sign methods stop ignoring rare events in heavy-tailed data (\cref{ex:sign:rare}).

\section{Back in the wild}\label{sec:sign:back}

\fieldnotes{Adam, Lion and signSGD move every parameter by about the learning rate per step: Lion and signSGD exactly, Adam within the factor $C(\beta_1,\beta_2)$.}
{Sign descent as normalized steepest descent in $\ell_\infty$, with $\nabla^2\Loss\preceq\diag(\vec L)$ and minibatch noise of variance $\sigma_i^2/B$.}
{Decrease $\eta\norm{\nabla\Loss}_1-\frac12\eta^2\norm{\vec L}_1$, exact on separable quadratics; rotations move $\norm{H}_{\infty\to1}$ between $\tr H$ and $d\lmax$; a sharp bias $\norm{\vec\sigma}_1/\sqrt B$, absent for symmetric unimodal noise; $\abs{m_t}/\sqrt{v_t}\le C(\beta_1,\beta_2)$, $7.27$ at the defaults, $1$ if $\beta_1=\beta_2$.}
{Sign descent wins when the gradient is spread over more coordinates than the curvature has stiff directions; a random rotation costs about $\frac2\pi d\lmax/\tr H$; signSGD ignores classes rarer than about $\ln2/B$; sign methods follow the square-root scaling rule (derived).}
{The memory that makes Adam and Lion work; why the coordinates of networks are good ones; the size of the Adam--SGD gap.}
{Use sign updates with momentum, not without, at small batch sizes (proved in the model); expect the advantage of Adam-like methods to depend on the basis (proved for quadratics, observed for GPT-2); a smaller $\beta_2$ keeps Adam's step near the learning rate (proved bound).}

What the theory tells a practitioner is modest and specific. A sign method with momentum suits problems whose gradients are spread over many parameters and whose curvature is concentrated in a few directions that the coordinates respect; whether a given network is such a problem the theory does not say. A sign method without momentum should not be trusted at small batch sizes, because it follows the median of the gradient, and the median ignores rare events. Adam's learning rate is the length of a step to within the factor $C(\beta_1,\beta_2)$, independent of the scale of the gradients when $\epsilon$ is negligible, which is why Lion, whose step always has RMS norm $1$, needs a smaller learning rate than Adam, whose update RMS is usually well below $1$. Claims beyond these, such as that Lion generalizes better than AdamW, are folklore as far as this chapter goes.

The spectral norm plays the role of the max norm for matrices (\cref{ch:spectral}); for a layer, sign descent moves the output by an amount proportional to the fan-in, which is why $\mu$P scales Adam's learning rate by $1/\mathrm{fan\_in}$ (\cref{ch:modular}). Momentum repairs the bias of the sign (\cref{ch:momentum}), the second moment turns the sign into a preconditioner (\cref{ch:adaptive}), the dependence on the basis returns with Shampoo and SOAP and with rotated transformers and rare tokens (\cref{ch:kronecker,ch:transformers}), and $C(\beta_1,\beta_2)$ bounds Adam's largest steps when an instability begins (\cref{ch:spikes}).

\begin{exercises}
\item \consolidation (Bias correction and equal coefficients.)\label{ex:sign:biascorr} Let $\rho=\beta_1^2/\beta_2<1$ and $\epsilon=0$. (a) Show that with bias correction $\abs{\hat m_t}/\sqrt{\hat v_t}\le B_t$, where
\[
  B_t=\frac{(1-\beta_1)\sqrt{1-\beta_2^{t+1}}}{(1-\beta_1^{t+1})\sqrt{1-\beta_2}}\Bigl(\frac{1-\rho^{t+1}}{1-\rho}\Bigr)^{1/2},\qquad B_0=1 .
\]
(b) Show that $B_t\le C(\beta_1,\beta_2)$ for every $t$, with equality for every $t$ if and only if $\beta_1=\beta_2$. (c) For $(0.9,0.999)$, compute $B_{10}$ and $B_{100}$ (answers: $1.05$ and $2.25$). (d) Let $\beta_1=\beta_2\in(0,1)$. Show that $\hat s_t^2=\hat v_t-\hat m_t^2$ is a weighted variance of the past gradients and that $\hat m_t/\sqrt{\hat v_t}=\sign(\hat m_t)(1+\hat s_t^2/\hat m_t^2)^{-1/2}$, so that \eqref{eq:sign:snr} is an identity in empirical moments; deduce $\abs{\hat m_t}\le\sqrt{\hat v_t}$, with equality if and only if $g_0=\dots=g_t$. \hint{with $a=\beta_1^{t+1}$ and $b=\beta_2^{t+1}$, the claim $B_t\le C$ is $(1-b)(1-a^2/b)\le(1-a)^2$.}

\item \consolidation\hard (Two-cycles of sign descent.)\label{ex:sign:cycles} Let $\Loss(\theta)=\frac12(\theta-\theta^\star)^\top H(\theta-\theta^\star)$ with $H\succ0$, and let $s\in\{\pm1\}^d$. (a) Show that $\{\theta,\theta-\eta s\}$ is an orbit of period two of sign descent if and only if $\sign(Hs)=s$ and the center $\bar\theta=\theta-\frac12\eta s$ satisfies $\abs{(H(\bar\theta-\theta^\star))_i}<\frac12\eta\abs{(Hs)_i}$ for every $i$. (b) Show that over the cycle $\norm{\nabla\Loss}_1$ averages $\frac12\eta s^\top Hs$ and the loss averages $\frac12(\bar\theta-\theta^\star)^\top H(\bar\theta-\theta^\star)+\frac18\eta^2s^\top Hs$. (c) For diagonal $H$, show that every two-cycle has average loss below $\frac14\eta^2\tr H$; in general, show that the centers form $\theta^\star+H^{-1}W$ for a box $W$, and explain the trap of \cref{fig:sign:rotation}. (d) Show that $s\mapsto\sign(Hs)$ (with $\sign0$ replaced by $1$) does not decrease $s^\top Hs$, and that a maximizer of $s^\top Hs$ satisfies $\sign(Hs)=s$. \hint{in (d), $(s'-s)^\top H(s'-s)\ge0$ for $s'=\sign(Hs)$, and flipping $s_i$ changes $s^\top Hs$ by $4H_{ii}-4s_i(Hs)_i$.}

\item \consolidation (Cantelli is sharp.) For $\mu,s,\delta>0$ let $X=\mu+\zeta$, where $\zeta=-\mu-\delta$ with probability $p=s^2/(s^2+(\mu+\delta)^2)$ and $\zeta=p(\mu+\delta)/(1-p)$ otherwise. Check that $\zeta$ has mean $0$ and variance $s^2$, compute $\mu-\mu\E\sign X$, and show that its supremum over $\mu$ and $\delta$ is $s$: the constant of \cref{lem:sign:signerror}(ii) cannot be lowered.

\item \consolidation\harder (Coordinatewise against $\ell_\infty$ smoothness.)\label{ex:sign:sdp} Let $H\succeq0$. (a) Show that $\min\{\norm{\vec L}_1:\diag(\vec L)\succeq H\}=\max\{\tr(HX):X\succeq0,\ X_{ii}=1\}$. (b) For feasible $X=VV^\top$, $\gamma\sim\cN(0,I)$ and $s=\sign(V\gamma)$, show that $\E[ss^\top]=\frac2\pi\arcsin(X)$ entrywise and that $\arcsin(X)\succeq X$. (c) Deduce $\norm{H}_{\infty\to1}\ge\frac2\pi\tr(HX)$ and the factor $\pi/2$ of \cref{rem:sign:sdp} \parencite{Nesterov1998}. \hint{(a) is semidefinite duality with strictly feasible points on both sides; in (b), $\E[\sign(a^\top\gamma)\sign(b^\top\gamma)]=1-2\angle(a,b)/\pi$, and the series of $\arcsin$ has nonnegative coefficients, so the Schur product theorem applies.}

\item \variation (Majority vote.)\label{ex:sign:vote} Let $M$ be odd, and let $M$ workers compute independent minibatch gradients whose coordinate noises are symmetric and unimodal with variances at most $\sigma_i^2/B$; the server steps by $-\eta\sign\bigl(\sum_m\sign(g^{(m)}_{t})\bigr)$. (a) With \cref{lem:sign:signerror}(iii) and Hoeffding's inequality, show that the vote has the wrong sign in coordinate $i$ with probability at most $\exp\bigl(-\frac M2\min\{S_i^2/3,\frac49\}\bigr)$, where $S_i=\sqrt B\abs{\partial_i\Loss}/\sigma_i$. (b) Prove the analogue of \cref{thm:sign:signsgd}. (c) Compare with one worker and batch $MB$ under Gaussian noise, whose error probability is $\Phi(-\sqrt MS_i)\le\frac12e^{-MS_i^2/2}$.

\item \variation (Lion.)\label{ex:sign:lion} (a) Show that Lion with $\beta_1=\beta_2$ is Signum. (b) Show that $c_t=\sum_{k=0}^tw_kg_{t-k}$ with $w_0=1-\beta_1$ and $w_k=\beta_1(1-\beta_2)\beta_2^{k-1}$ for $k\ge1$, that the weights sum to $1-\beta_1\beta_2^t$, and that their mean lag $\sum_kkw_k$ tends to $\beta_1/(1-\beta_2)$, which is $90$ steps for the defaults; compare with Signum at $\beta=\beta_2$. (c) Lion's update has RMS norm $1$ when no coordinate of $c_t$ vanishes. Using \eqref{eq:sign:snr}, explain why matching AdamW's movement per step calls for a Lion learning rate several times smaller.

\item \variation (Memoryless Adam with $\epsilon$.)\label{ex:sign:epsilon} Let $\theta_{t+1}=\theta_t-\eta\,\psi(\nabla\Loss(\theta_t))$ with $\psi(x)=x/(\abs x+\epsilon)$ coordinatewise. (a) Under \cref{ass:sign:smooth}, show that $\Loss(\theta_{t+1})\le\Loss(\theta_t)-\eta\sum_ig_i^2/(\abs{g_i}+\epsilon)+\frac12\eta^2\norm{\vec L}_1$ for $g=\nabla\Loss(\theta_t)$. (b) Show that $x^2/(\abs x+\epsilon)\ge\frac12\min\{\abs x,x^2/\epsilon\}$, and deduce a bound of the form \eqref{eq:sign:mixed} with $\epsilon$ in the role of $\sigma_i/\sqrt B$. (c) When does $\epsilon=10^{-8}$ matter?

\item \variation (Block-diagonal curvature.) Let $H=\diag(H_1,\dots,H_K)$ with blocks of sizes $d_1,\dots,d_K$. (a) Show that $\norm{H}_{\infty\to1}=\sum_k\norm{H_k}_{\infty\to1}$. (b) If each block is rotated by its own orthogonal matrix, show that the $\ell_\infty$ constant is at most $\sum_kd_k\lmax(H_k)$, and compare with $d\lmax(H)$ for a rotation that mixes the blocks. (c) Relate this to the block structure that \textcite{ZhangEtAl2024} measure in transformers.

\item \simulation (Shapes of noise.) With \texttt{02-sign-noise.py}: (a) vary $p$ and $B$ and check that the rare-event floor at small $\eta$ collapses as soon as $(1-p)^B<\frac12$. (b) Replace the rare-event noise by $\sigma(E-1)$ with $E$ exponential of mean $1$. Show that as $\eta\to0$ the floor at $B=1$ is the median gap $\sigma(1-\ln2)=0.31\sigma$, and that for $B\ge2$ it is within $3\%$ of $\sigma/(3B)$, the median gap of a gamma law. Conclude that for a fixed skewed law the bias decays like $1/B$, and that the worst case $1/\sqrt B$ of \cref{thm:sign:signsgd} needs events of probability comparable to $1/B$.

\item \simulation (Adam after a rotation.)\label{ex:sign:adamrot} With \texttt{02-sign-rotation.py}, run Adam with $(\beta_1,\beta_2)=(0.9,0.999)$, $\epsilon=0$ and $\eta=0.01$ on both problems. You should find a loss near $10^{-9}$ after $1000$ steps on the diagonal problem and near $3\times10^{-3}$ on the rotated one, still ahead of gradient descent. Show that along the rotated run the preconditioned sharpness $\eta\lmax(P^{1/2}Q^\top HQP^{1/2})$, with $P=\diag(\hat v_t)^{-1/2}$, grows towards the stability threshold $(2+2\beta_1)/(1-\beta_1)=38$ of momentum (\cref{ch:momentum,ch:eos}), and explain why Adam's advantage there has nothing to do with the coordinates.

\item \research (When are the coordinates good?)\label{ex:sign:basis} (a) Show that $\norm{H}_{\infty\to1}\le\sum_j\lambda_j\norm{v_j}_1^2$ for $H=\sum_j\lambda_jv_jv_j^\top$, and find conditions on the localization of the eigenvectors under which $\norm{H}_{\infty\to1}/\tr H=O(1)$. (b) Estimate both sides for the Hessian of a small network with Hessian--vector products and the iteration of \cref{ex:sign:cycles}(d). What is open is a mechanism that explains why the Hessians of trained networks are closer to diagonal, or block diagonal, in the parameters' coordinates than a random rotation would make them.

\item \research (Rare events and memory.)\label{ex:sign:rare} (a) For Signum with momentum $\beta$ in the model of \cref{prop:sign:rareclass}, argue that the update follows the sign of an average of about $B(1+\beta)/(1-\beta)$ labels, derive the class frequency below which Signum fails, and check it by simulation. (b) For a softmax output layer with Zipf frequencies $\pi_k\propto1/k$ over $K$ classes, find the fraction of classes that signSGD and Signum learn as a function of $B$ and $\beta$. What is open is whether this mechanism affects the rare tokens of language models trained with Adam, whose $\beta_2<1$ restores the mean (\cref{ch:adaptive,ch:transformers}).
\end{exercises}

\notesandrefs

Sign methods are older than their analysis (\cref{sec:sign:wild}). \Textcite{BernsteinEtAl2018} gave the first nonconvex analysis of signSGD and Signum under coordinatewise smoothness, with the comparison of the densities of gradients and curvatures that \cref{cor:sign:compare} restates, and \textcite{BernsteinEtAl2019} analysed majority vote under symmetric unimodal noise. \Textcite{KarimireddyEtAl2019} showed that signSGD can fail to converge and that error feedback repairs it; \textcite{SafaryanRichtarik2021} replaced the noise assumptions by a condition on the probability that each sign is correct. We know of no source for \cref{prop:sign:floor,prop:sign:rareclass} in this form; both are elementary, and the second is the median-against-mean mechanism of the counterexample of \textcite{ReddiKaleKumar2018}.

Steepest descent in the max norm is a tool outside machine learning as well; \textcite{KelnerEtAl2014} used it in an almost-linear-time algorithm for approximate maximum flow. \Textcite{BallesPedregosaLeRoux2020} compared coordinatewise and $\ell_\infty$ smoothness for sign descent and related the latter to the concentration of the Hessian on its diagonal; the factor $\pi/2$ of \cref{rem:sign:sdp} is Nesterov's theorem on the semidefinite relaxation of quadratic maximization over the cube \parencite{Nesterov1998}. \Textcite{BernsteinNewhouse2024} organized optimizers by norms and identified Adam without its averages as steepest descent in $\ell_\infty$, and \textcite{XieMohamadiLi2025} developed the $\ell_\infty$ theory of Adam and tested it by rotating GPT-2's parameters.

\Textcite{KingmaBa2015} introduced Adam and discussed the size of its effective step; the supremum $C(\beta_1,\beta_2)$ takes one application of the Cauchy--Schwarz inequality. \Textcite{BallesHennig2018} separated Adam's sign from its variance adaptation, \textcite{KunstnerEtAl2023} found that Adam behaves like sign descent with momentum as the batch grows, and \textcite{ZhaoEtAl2025} that Signum matches Adam on language models. \Textcite{MalladiEtAl2022} derived the square-root scaling rule for adaptive methods. Lion was found by \textcite{ChenEtAl2023} through a symbolic program search; \textcite{ChenLiuLiangLiu2024} showed with a Lyapunov function that Lion with weight decay solves a box-constrained problem, \textcite{XieLi2024} proved the corresponding statement for AdamW, and \textcite{ZhangZouCao2024} that Adam on separable data converges in direction to an $\ell_\infty$ max-margin classifier.

\paragraph{Translating the literature.} \Textcite{BernsteinEtAl2018} write $f$ for our $\Loss$, $x_k$ for $\theta_t$, $g_k$ for the full gradient $\nabla\Loss(\theta_t)$ and $\tilde g_k$ for the stochastic gradient $g_t$, $\delta_k$ for the step size $\eta_t$, $n_k$ for the batch size $B$ and $K$ for the number of steps $T$; their $\vec L$ and $\vec\sigma$ are ours. \Textcite{KingmaBa2015} write $\alpha$ for the step size and start at $t=1$ with $m_0=v_0=0$, so their $\hat m_t=m_t/(1-\beta_1^t)$ is our $\hat m_{t-1}$. \Textcite{ChenEtAl2023} write Lion with $c_t$ for the signed interpolation, $m_t$ for the momentum and $\lambda$ for the decoupled weight decay, as we do. Some implementations of the sign map $0$ to $1$; in this chapter that changes events of probability zero only.

\printchapterbib
\end{refsection}
